\section{Introduction}
\label{sec:introduction}

Let \(S\) be the sphere spectrum, fix a prime \(p\), and write
\(S_p^\wedge\) for its \(p\)-completion. We use \(L_i\) for localization
at Johnson--Wilson theory \(E(i)\) and \(L_{K(i)}\) for localization at
Morava \(K\)-theory of height \(i\), all at \(p\). In particular, \(L_0\)
is rationalization. Chromatic splitting concerns the lower-height
overlap \(L_{n-1}L_{K(n)}S\) in the fracture square for \(L_nS\).
Hopkins' strong formulation, recorded by Hovey
\cite[Conjecture~4.2]{Hovey1995}, prescribes an assembly of this overlap
from localized spheres.

At height three, the proposed underlying spectrum is
\begin{equation}\label{eq:proposed-wedge}
\begin{aligned}
 \mathcal W_3={}&L_2\bigl(S_p^\wedge\vee\Sigma^{-1}S_p^\wedge\bigr)\\
 &{}\vee L_1\bigl(\Sigma^{-3}S_p^\wedge\vee\Sigma^{-4}S_p^\wedge\bigr)\\
 &{}\vee L_0\bigl(\Sigma^{-5}S_p^\wedge\vee\Sigma^{-6}S_p^\wedge
             \vee\Sigma^{-8}S_p^\wedge\vee\Sigma^{-9}S_p^\wedge\bigr).
\end{aligned}
\end{equation}
The strong prediction is an equivalence
\(L_2L_{K(3)}S\simeq\mathcal W_3\), with specified maps, including the
canonical unit on the first summand. The revised formulation of
Barthel--Beaudry retains this odd-prime prediction
\cite[Conjecture~6.3 and Remark~6.4]{BB2019}. We disprove even the
underlying equivalence of \(E(2)\)-local spectra.

The obstruction is a natural map. For every spectrum \(X\), the
\(K(2)\)-localization unit induces
\[
 \tau_X:L_0X\longrightarrow L_0L_{K(2)}X.
\]

\begin{theorem}\label{thm:main}
For every prime \(p\geq5\), the canonical map
\begin{equation}\label{eq:main-map}
 \tau_{L_{K(3)}S}:L_0L_{K(3)}S\longrightarrow
                         L_0L_{K(2)}L_{K(3)}S
\end{equation}
is nonzero on \(\pi_{-3}\).
\end{theorem}

\begin{corollary}\label{cor:strong}
For every prime \(p\geq5\), there is no equivalence of \(E(2)\)-local
spectra \(L_2L_{K(3)}S\simeq\mathcal W_3\), with \(\mathcal W_3\) as in
\eqref{eq:proposed-wedge}. Thus the height-three strong chromatic
splitting formula is false in this range, whether or not an equivalence
is required to preserve its specified unit summand.
\end{corollary}

\subsection{Why the canonical map matters}

The rational homotopy groups of the \(K(n)\)-local sphere are known
for every positive height \(n\) and every prime. Barthel--Schlank--Stapleton--Weinstein
identify their algebra with an exterior algebra over \(\Q_p\), with
primitive generators in degrees \(1-2i\), \(1\leq i\leq n\)
\cite[Theorem~A]{BSSW2025}. At height three the resulting degrees are
\[
 0,-1,-3,-4,-5,-6,-8,-9,
\]
exactly the shifts appearing in \eqref{eq:proposed-wedge}. Thus the
rational degree pattern does not distinguish the proposed wedge from
the actual overlap.

The natural transformation \(\tau\) does distinguish them. In
\(\mathcal W_3\), a degree-\(-3\) class can occur only in the
lower-height part, which is \(K(2)\)-acyclic. The two top-height
summands have no rational homotopy in that degree. It follows that
\(\tau_{\mathcal W_3}\) vanishes on \(\pi_{-3}\). The formal argument
in \cref{sec:wedge-obstruction} also compares the maps for
\(L_{K(3)}S\) and \(L_2L_{K(3)}S\). Naturality makes the conclusion
valid for any hypothetical equivalence, with no assumption on its
individual summand maps.

At height two, Hovey records Hopkins' splitting
for \(p>3\), based on Shimomura--Yabe's calculations
\cite[\S4]{Hovey1995}, \cite[Theorem~2.4]{ShimomuraYabe1995}. Behrens later reorganized those calculations
and displayed the predicted overlap homotopy groups
\cite[Remark~7.8]{Behrens2012}. Goerss--Henn--Mahowald proved the
splitting at \(p=3\) \cite[Theorem~5.11]{GHM2014}. At \(p=2\),
Beaudry disproved the original formula \cite[Theorem~1.4]{Beaudry2017},
while Beaudry--Goerss--Henn subsequently proved a refined splitting
with additional Moore-spectrum summands and retained the weak unit
splitting \cite[Theorem~1.1.6]{BGH2022}.

The canonical map already controls the height-two argument of
Goerss--Henn--Mahowald. At \(p=3\), their proof identifies the rational
map from the \(K(2)\)-local sphere to its \(K(1)\)-localization with
projection onto the degrees zero and minus one, followed by canonical
localization; the fracture square then reconstructs the splitting
\cite[proof of Theorem~5.11, p.~1288]{GHM2014}.
Theorem~\ref{thm:main} finds a nonzero degree-minus-three component at
height three, which obstructs \(\mathcal W_3\). A finite filtration
with these cofibers may retain nonzero attaching maps that a wedge
would discard. The existence of a retraction of the canonical unit
in weak chromatic splitting remains a separate question.

\subsection{A detecting spectrum}

The proof of Theorem~\ref{thm:main} has two tasks: construct a map
from \(L_{K(3)}S\) to an ordinary \(K(2)\)-local spectrum, and prove
that this map preserves a nonzero rational degree-\(-3\) class. Locality
then forces the map to factor through the canonical localization unit.
The class is detected in degree-three stabilizer cohomology; the main
argument compares its height-three and height-two realizations.

Fix \(p\geq5\), put \(k=\F_{p^6}\), and choose \(C=\C_p\) with
compatible constant-field data. For \(n=2,3\), let \(\Gamma_n\) be the
height-\(n\) Honda formal group over \(k\), and put
\[
 E_n=E(k,\Gamma_n),\qquad
 P_n=\Aut_k(\Gamma_n)=\cO_{D_n}^{\times}.
\]
Here \(D_n\) is the division algebra over \(\Q_p\) of invariant
\(1/n\), and \(P_n\) is the nonextended stabilizer. Set
\[
 G=P_3\times P_2,\qquad
 \Delta=\Gal(k/\F_p),\qquad G_3=P_3\rtimes\Delta.
\]
The field \(k\) contains the endomorphism fields of both Honda groups.

On the height-exactly-two stratum \(k((u_2))\) of the height-three
deformation space, the Gross--Torii extension identifies the connected
height-two group with \(\Gamma_2\). Its valued-field completion
\(\widehat L\) carries the two framing actions, hence an action of
\(G\). The solid Lubin--Tate construction of
Barthel--Mann--Ray--Schlank--Senger--Weinstein--Zhou \cite{BMR2026}
gives the connected completed theory
\[
 \widehat B^\square=E^\square(\widehat L,(\Gamma_2)_{\widehat L}),
 \qquad \mathcal D=\bigl((\widehat B^\square)^{hG}\bigr)(*).
\]
The superscript \(\square\) records the solid structure and continuous
action; \((*)\) is evaluation in ordinary spectra. Fixed points are
formed before this evaluation.

Section~\ref{sec:completed-tower} constructs equivariant maps
\(E_3^\square\to\widehat B^\square\leftarrow E_2^\square\).
The first requires a section on oriented deformations, obtained by
adjoining the marked \'etale factor to the connected deformation.
The second is the canonical connected base-change map. Restriction
from \(G_3\) to \(P_3\), inflation along \(G\to P_3\), and the first
map define \(g:L_{K(3)}S\to\mathcal D\). Every profinite evaluation
of \(\widehat B^\square\) is \(K(2)\)-local, so its bar totalization
\(\mathcal D\) is too. Thus
\begin{equation}\label{eq:overview-detector}
 L_{K(3)}S\xrightarrow{\eta}L_{K(2)}L_{K(3)}S
       \xrightarrow{\overline g}\mathcal D,
 \qquad g=\overline g\eta.
\end{equation}
The source identification and these actual maps are established in
\cref{prop:completed-theory,lem:local-detector}.

Write \(N_3:P_3\to\Z_p^\times\) for reduced norm and
\(P_3^1=N_3^{-1}(\mu_{p-1})\). The completed comparison
\[
 E_2^\square\xrightarrow{\sim}(\widehat B^\square)^{hP_3^1}
\]
reduces the detector's coefficients to height two. We prove the
individual-coefficient comparison as well as this instance of
\cite[Theorem~5.2.6]{BMR2026}. The corresponding equivalence for the
uncompleted half sphere remains a separate conjecture
\cite[Conjectures~1.1.1 and~5.1.13]{BMR2026}, with a conditional
coheight-one consequence in \cite[Proposition~5.1.17]{BMR2026}.
The completed target in \eqref{eq:overview-detector} suffices here.

\subsection{Transporting the primitive}

The starting classes are nonzero elements
\(e_{P_n}\in H^3_{\cts}(P_n,\Q_p)\), for \(n=2,3\).
Under rational Lie comparison they are represented, up to nonzero
scalars, by
\[
 (U,V,W)\longmapsto\tr_{\mathrm{red}}\bigl(U[V,W]\bigr).
\]
This is the classical cocycle attached to an invariant bilinear form
\cite[\S21]{ChevalleyEilenberg1948}, \cite[\S11]{Koszul1950}.
Section~\ref{sec:cohomology} constructs these classes and their linear
counterparts \(e_{J_n}\) for
\[
 J_n=\{g\in\GL_n(\Q_p):|\det g|_p=1\}.
\]
The trace class restricts nontrivially from rank three to rank two;
the same compatibility holds for the relevant parabolic subgroup,
with determinant lattices retained. The low-rank calculation and the building argument are given
in \cref{app:background}.

To compare division-algebra and linear classes, we use a line
modification of a form of the rank-\(n\), degree-one Fargues--Fontaine
bundle \(\mathcal E_n^0=\cO(1/n)\), retaining a determinant lattice.
The stack \(\mathcal M_n\) of such modifications has two presentations:
\[
 \mathcal M_n\simeq[\mathbb P_C^{n-1,\diamondsuit}/P_n]
       \simeq[\Omega_C^{n-1,\diamondsuit}/J_n].
\]
Here projective space parametrizes quotient lines of the untilt fiber,
and Drinfeld space \(\Omega_C^{n-1}\) is the complement of the
\(\Q_p\)-rational hyperplanes. The relative bundle correspondence and
basic-stratum theorems of Fargues--Scholze provide the torsors in
families \cite[Theorem~II.2.19, Corollary~II.2.20, and
Theorems~III.2.4 and~III.4.5]{FS}. Keeping a lattice in the determinant
local system gives the structure group \(J_n\). These presentations
follow Faltings's Lubin--Tate/Drinfeld duality
\cite[Introduction]{Faltings2002}, in the perfectoid form of
Scholze--Weinstein with its exchanged period maps
\cite[Theorem~7.2.3]{SW2013}. Section~\ref{sec:modifications}
constructs the determinant-lattice comparison used here.

For the geometric cohomology below, write \(H^i_{\mathrm b}\) for
integral cohomology followed by inversion of \(p\). Integral cohomology
is the derived inverse limit of finite \(p\)-power coefficient complexes,
as specified in Section~\ref{sec:cohomology}.
The two presentations have complementary roles. The projective bundle
formula singles out a one-dimensional Galois-invariant line in
\(H^3_{\mathrm b}(\mathcal M_n)\), containing the nonzero image of \(e_{P_n}\).
Using the Galois weights in the integral Drinfeld cohomology theorem
of Colmez--Dospinescu--Nizioł \cite[Theorems~1.1 and~5.1]{CDN2021},
we prove that the image of \(e_{J_n}\) is also nonzero and belongs
to that line. Thus the two
classes are proportional.

The completed tower carries a surjection from the rank-three bundle
to the rank-two bundle. A common line modification turns this into
an exact sequence of rational local systems of ranks \(1,3,2\).
Parabolic restriction and projective-bundle injectivity then give
\[
 e_{P_3}|_{Y_C}=a\,e_{P_2}|_{Y_C},\qquad a\in\Q_p^\times,
 \quad
 Y_C=[\Spd\widehat L/G]\times_{\Spd k}\Spd C.
\]
Section~\ref{sec:degree-three} proves this relation. Both classes on
\(Y_C\) could still vanish, so nonvanishing requires a further step.

Section~\ref{sec:witt-transfer} first reflects the relation from
\(Y_C\) to \(Y=[\Spd\widehat L/G]\) with Witt coefficients. The
maps from constants become equivalences after \(P_3^1\)-descent;
a bounded finite free complex proves injectivity under the change
of constants \(W(k)\to W(C^\flat)\). The continuous additive
retraction onto Witt constants of Barthel--Schlank--Stapleton--Weinstein
\cite[Proposition~2.5.1 and Corollary~2.5.9]{BSSW2025} then proves
nonvanishing of the height-two class in deformation coefficients.
If \(c_n\) is the image of \(e_{P_n}\) under projection from \(G\)
and the coefficient unit, the result is
\[
 c_3=a\,c_2\ne0
 \quad\text{in }H^3(G,\pi_0\widehat B^\square)[1/p].
\]

Finally, Section~\ref{sec:detection} uses the descent spectral
sequences for the actual map \(g\). A finite cohomological bound
produces a finite integral filtration in each ordinary homotopy
degree. After rationalization, central scalar automorphisms kill
all coefficient degrees \(t\ne0\). With \(s\) denoting group-cohomological
degree, the nonzero coefficient map in bidegree \((s,t)=(3,0)\) therefore survives to degree \(-3\) ordinary
homotopy. The factorization \eqref{eq:overview-detector} proves the
main theorem, and Section~\ref{sec:wedge-obstruction} gives its
wedge consequence.

Continuous invariants throughout the proof retain their solid
coefficient objects, and inversion of \(p\) follows integral cohomology.
This order matters for the noncompact spaces and groups that occur here.
